\originalchapter{8.512作业1--3：涨落、介电与Peierls不稳定性}
\psource{Patrick A. Lee，8.512 Spring 2009，公开作业1--3；题页保留2008底稿和2009改期痕迹，不把印刷日期混作独立课次。}
\originalsection{作业1（公开修订题页，2009年2月）}
\begin{exercise}[原题1(a)--(b)]\label{p:512-1-1}
原题取$\hbar=1$，$S(q,\omega)=\int\dd x\dd t\,\e^{-\ii qx+\ii\omega t}\avg{\rho(x,t)\rho(0,0)}$，$n_B=(\e^{\beta\omega}-1)^{-1}$，响应采用外势$+f\rho$约定。(a)证明$\chi''=(\e^{-\beta\omega}-1)S/2$与$S=-2(n_B+1)\chi''$。(b)证明$\chi''(q,\omega)=-\chi''(-q,-\omega)$、$S(-q,-\omega)=\e^{-\beta\omega}S(q,\omega)$，解释详细平衡。
\end{exercise}
\begin{solution}
插入能量态，$S=2\pi\sum_{nm}p_n|\rho_{nm}(q)|^2\delta(\omega-E_m+E_n)$。交换$n,m$使$q,\omega$均反转，权重比$p_m/p_n=\e^{-\beta\omega}$，得到详细平衡。该响应约定的谱虚部为$-\pi\sum(p_n-p_m)|\rho_{nm}|^2\delta(\omega-E_m+E_n)$，故$\chi''=-(1-\e^{-\beta\omega})S/2$。利用$n_B+1=(1-\e^{-\beta\omega})^{-1}$得第二式。两次反转令虚部反号。探针失能$\omega>0$的概率相对于得能大$\e^{\beta\omega}$倍，是样品吸收与释放能量的热平衡比例。
\end{solution}
\begin{exercise}[原题2(a)--(e)]\label{p:512-1-2}
晶体核位置$r_j=R_j+u_j$，$\hbar=1$，$u_j=\sum_q(2NM\omega_q)^{-1/2}(a_q\e^{\ii(qR_j-\omega_qt)}+\mathrm{h.c.})$；先假设三种极化简并，选择沿$Q$的一种并将$Q\cdot u_j$写为$Qu_j$。定义$F(Q,t)=\sum_{jl}\e^{-\ii Q(R_j-R_l)}F_{jl}(t)$，$F_{jl}=\avg{\e^{-\ii Qu_j(t)}\e^{\ii Qu_l(0)}}$，$S=\int\e^{\ii\omega t}F\dd t$。(a)证明$F_{jl}=\avg{\e^{-\ii Q[u_j(t)-u_l(0)]}}\e^{[Qu_j(t),Qu_l(0)]/2}$，可调用高斯平均公式。(b)令$\theta_{jl}=q(R_j-R_l)-\omega_qt$，证明
\begin{equation}
F_{jl}=\e^{-2W}\exp\left\{\frac{Q^2}{2NM}\sum_q\frac{(2n_q+1)\cos\theta_{jl}+\ii\sin\theta_{jl}}{\omega_q}\right\},\quad2W=\frac{Q^2}{2NM}\sum_q\frac{2n_q+1}{\omega_q}.
\end{equation}
(c)展开指数至一个声子阶，写出弹性及$Q=G\pm q,\omega=\pm\omega_q$项，$V^*$为倒格原胞体积。(d)解释各项。(e)怎样测声子极化向量？
\end{exercise}
\begin{solution}
(a)谐振位移的不同时间对易子为数，BCH恰截断，因$(-\ii)(\ii)=1$得到所给正号。高斯平均为$\exp[-Q^2\avg{(u_j-u_l)^2}/2]$。(b)位移自相关的两个本征权重为$(n_q+1)\e^{\ii\theta}+n_q\e^{-\ii\theta}=(2n_q+1)\cos\theta+\ii\sin\theta$；自方差项组成$2W$。

(c)把该恒等式代入指数的一阶，格点和给倒格$\delta$，时间积分给能量$\delta$。若用通常$\int\e^{\ii\omega t}\dd t=2\pi\delta(\omega)$，完整表达为
\begin{align}
S=2\pi NV^*\e^{-2W}\bigg[&\sum_G\delta(Q-G)\delta(\omega)\\
+&\frac{Q^2}{2NM}\sum_q\frac1{\omega_q}\Big\{(n_q+1)\sum_G\delta(Q-G-q)\delta(\omega-\omega_q)\nonumber\\
&\hspace{31mm}+n_q\sum_G\delta(Q-G+q)\delta(\omega+\omega_q)\Big\}\bigg].\nonumber
\end{align}
原题把Fourier归一的$2\pi$吸收到其$\delta$记号中；须同时保持$S$定义与实验散射率的约定。(d)Bragg弹性峰、创造一个声子、吸收一个声子依次对应三项；零温吸收项为零。(e)恢复非简并极化后强度为$|Q\cdot\lambda_{q\alpha}|^2$，改变$Q$方向并在相同$q,\omega_{q\alpha}$测强度可重建投影，纵横模需几何与多方向数据区分。
\end{solution}
\originalsection{作业2（题页修订为2009年2月23日）}
\begin{exercise}[原题1(a)--(c)]\label{p:512-2-1}
高速电子在金属内因发射等离子体激元而失能。(a)小$q$时$\operatorname{Im}[-1/\epsilon(q,\omega)]\simeq A(q)\delta(\omega-\omega_p)$，由和规则求$A$。(b)写出发射一个激元进入固角的概率表达式。(c)取电子能100 keV等典型值估计平均自由程，说明选取参数。
\end{exercise}
\begin{solution}
(a)逆介电和规则$\int_0^\infty\omega\operatorname{Im}[-\epsilon^{-1}]\dd\omega=\pi\omega_p^2/2$，所以$A=\pi\omega_p/2$。同样可由$\epsilon=1-\omega_p^2/(\omega+\ii0)^2$直接取极点检验。

(b)选择平面波归一体积$V$与能量变量$E=\hbar\omega$，外电子的屏蔽库仑势谱给微分率
\begin{equation}
\dd\Gamma=\frac2\hbar\frac{\dd^3q}{(2\pi)^3}\,v_q\operatorname{Im}[-\epsilon^{-1}(q,\omega_q)]\,,\qquad \hbar\omega_q=\epsilon_k-\epsilon_{k-q},\quad v_q=\frac{4\pi e_*^2}{q^2}.
\end{equation}
代入$A\delta(\omega_q-\omega_p)$再由$q=k-k'$换到$k'^2\dd k'\dd\Omega$，并对$k'$积分，就得定角发射率；可写为$\dd\Gamma/\dd\Omega=(2/\hbar)(2\pi)^{-3}v_{|k-k'|}A\,k'^2/|\partial_{k'}\omega_q|$，$k'$由能量守恒确定。该表达指定了归一与Jacobian，避免把速度与波矢单位混在一起。

(c)在$v$远大于电子费米速度时忽略反冲，$\omega_q\simeq\vct v\cdot\vct q$，角积分给$1/(vq)$，于是$\Gamma\simeq(e_*^2\omega_p/\hbar v)\log(q_{\max}/q_{\min})$，$q_{\min}=\omega_p/v$，$\lambda=v/\Gamma$。取$\hbar\omega_p=15$ eV，100 keV电子的相对论速度$v=c\sqrt{1-(1+100/511)^{-2}}\simeq0.548c$；取激元进入粒子--空穴连续谱的截断$q_{\max}\sim\omega_p/v_F$、$v_F=1.5\times10^6$ m/s，得到对数约4.7、平均自由程约$1.15\times10^{-7}$ m量级。若采用其他截断，数字可变化；大$q$时单极点近似失效，这不是材料专属精确寿命。
\end{solution}
\begin{exercise}[原题2(a)--(c)]\label{p:512-2-2}
(a)证明周期固体的RPA介电函数为$1+v_q\sum_{k\alpha\beta}|\avg{k+q,\beta|\e^{\ii q\cdot r}|k,\alpha}|^2(f_{k\alpha}-f_{k+q,\beta})/(E_{k+q,\beta}-E_{k\alpha}-\hbar\omega-\ii0)$，各自旋与体积计数要一致。(b)静态长波半导体以平均带隙$\Delta$替代跃迁能；证明$\sum_b(E_b-E_a)|\avg{b|\e^{\ii q r}|a}|^2=\hbar^2q^2/(2m)$，说明所需平衡态条件。(c)再把该和中的跃迁能取$\Delta$，求$\epsilon(0)=1+(\hbar\omega_p/\Delta)^2$，估计Si、Ge并讨论与实验比较的限制。
\end{exercise}
\begin{solution}
(a)先在Bloch能量基对外势作一阶响应，矩阵元加占据差给粒子--空穴泡图，再自洽加Hartree势形成$\epsilon=1+v_q\Pi$，每个$k$应带$1/V$，自旋若不在$\alpha$中再乘2。

(b)$O=\e^{\ii q r}$满足$O^\dagger HO=H+\hbar\vct q\cdot\vct p/m+\hbar^2q^2/(2m)$。在无流的时间反演对称态或成对平均中$\avg{p}=0$，完备关系给所需式。对一个具有非零流的单Bloch态，还应保留$\hbar q\cdot\avg p/m$；原题提时间反演是为使该项取消，并不保证任意单个$k$态本身无流。

(c)绝缘体的带间贡献包含正反跃迁，对占据态求和后$\Pi\simeq n\hbar^2q^2/(m\Delta^2)$，故$\epsilon-1=\hbar^2\omega_p^2/\Delta^2$。若用价电子等离子体能约16 eV与有效平均跃迁尺度Si约4.9 eV、Ge约4.1 eV，则预测约11.7、16.2。这里选择平均光学跃迁尺度是模型输入，不能把间接最低带隙Si约1.1 eV、Ge约0.7 eV代入并宣称成功：那会给大得多的介电常数。与实验做精确比较需核对温度、频率及实测光谱；本题仅给模型量级及失效原因。
\end{solution}
\originalsection{作业3（题页修订为2009年3月2日）}
\begin{exercise}[原题1(a)--(e)]\label{p:512-3-1}
一维H链间距$a$，每原子一个$1s$轨道、一电子，只保留最近邻$V(a)<0$并忽略重叠。(a)求带与费米能。(b)原子位移$(-1)^n\delta$，令$\Delta=2\delta\dd V/\dd a$，证明跃迁交替$V\pm\Delta$并求新带、金属或绝缘性质。(c)求小$|\Delta/V|$时电子能改变的$\Delta^2\log|\Delta/V|$项及系数。(d)晶格代价每站$b\delta^2$，证明均匀链不稳定。(e)求一维自由费米气静态极化，证明$2k_F$对数奇异。
\end{exercise}
\begin{solution}
(a)$E(k)=\epsilon_{1s}+2V\cos ka$，双自旋半填充，$k_F=\pi/(2a)$、$E_F=\epsilon_{1s}$。(b)相邻距为$a\pm2\delta$，Taylor展开给交替跳跃。两站胞$q\in[-\pi/(2a),\pi/(2a)]$，$E_\pm=\epsilon_{1s}\pm2\sqrt{V^2\cos^2qa+\Delta^2\sin^2qa}$。下带填满，边界隙$4|\Delta|$，变成带绝缘体。

(c)按每原子并含自旋，电子能$E_{\mathrm{el}}/N=-(4/\pi)\int_0^{\pi/2}\sqrt{V^2\cos^2x+\Delta^2\sin^2x}\dd x$，忽略共同$\epsilon_{1s}$。边界$x=\pi/2-y$处差$\sqrt{V^2y^2+\Delta^2}-|V|y$，积分的对数部分$[\Delta^2/(2|V|)]\log(|V|/|\Delta|)$，所以$\delta E_{\mathrm{el}}/N=[2\Delta^2/(\pi|V|)]\log|\Delta/V|+O(\Delta^2)$。(d)$\Delta\propto\delta$，负$\delta^2\log(1/|\delta|)$压过任何固定正$b\delta^2$，在该零温无限链模型中任意小畸变可降能。有限温、有限长或电子关联改变红外截止。

(e)取双自旋态与$\Pi=-\chi$，一维积分直接得
\begin{equation}
\Pi(q,0)=\frac{2m}{\pi\hbar^2|q|}\log\left|\frac{|q|+2k_F}{|q|-2k_F}\right|.
\end{equation}
$q\to0$极限$2m/(\pi\hbar^2k_F)$，等于双自旋态密度；$q=2k_F$有对数，来源于两个费米点被同一波矢连接。
\end{solution}
\begin{exercise}[原题2]\label{p:512-3-2}
电子仅在$x,y$平面运动而电磁场在三维，求长波二维等离子体色散并证明$\omega_p\propto\sqrt{|q|}$。
\end{exercise}
\begin{solution}
平面库仑Fourier为$v_q=e^2/(2\epsilon_0\epsilon_r|q|)$；惯性密度响应高频为$\chi_0\simeq n_{2D}q^2/(m\omega^2)$，自洽条件$1-v_q\chi_0=0$给$\omega_p^2=n_{2D}e^2|q|/(2m\epsilon_0\epsilon_r)$。若错误把场也限成二维使势Fourier不同，就会改变色散。该式忽略迟滞、有限厚度及强Landau阻尼。
\end{solution}
